% █▀█ █▀█ █▀▀ ▄▀█ █▀▄▀█ █▄▄ █   █▀▀
% █▀▀ █▀▄ ██▄ █▀█ █ ▀ █ █▄█ █▄▄ ██▄

%::::::::::::::::::::::::::::::::::::::::::::::::::::::::::::::::::::::::::::::::
%                                   CORE SETUP                                   
%                             Document-Wide Settings                             
%................................................................................
\usepackage[
	a4paper,
	margin=2cm,
]{geometry}	% Set page margins
\usepackage[brazil]{babel}		% Language-specific settings
\usepackage{csquotes}			% Context-sensitive quoting
					% 'csquotes' works in tandem with babel
					% 'csquotes' IS A 'biblatex' DEPENDENCY

%::::::::::::::::::::::::::::::::::::::::::::::::::::::::::::::::::::::::::::::::
%                               MATH & FONT ENGINE                               
%................................................................................
% AMS (American Mathematical Society)
\usepackage{amsmath}		% AMS's package for mathematical typesetting
\usepackage{mathtools}		% Extends amsmath's functionalities (i.e.: a superset)
\usepackage{amsthm}		% AMS's package for creating theorem-like environments

\usepackage{fontspec}		% Interface for using system fonts (OTF/TTF) with Lua/XeLaTeX
\usepackage[
	math-style = ISO,
	bold-style = ISO
]{unicode-math}	% Must be loaded after 'mathtools'

% Font selection
\setmainfont{Libertinus Serif}
\setsansfont{Libertinus Sans}[
	BoldItalicFont = {Libertinus Sans Italic} % Fallback for missing BI
]
\setmonofont{Libertinus Mono}[
	BoldFont		= {Libertinus Mono},	% Fallback for missing B
	ItalicFont		= {Libertinus Mono},	% Fallback for missing I
	BoldItalicFont	= {Libertinus Mono}	% Fallback for missing BI
]
\setmathfont{Libertinus Math}
\setmathfont[range=\mathcal, StylisticSet=1]{XITS Math}	% Libertinus uses a single font for both script and calligraphic text
\setoperatorfont{\mathsf}	% Sans (without) serif font to distinguish mathematical operators from prose
%::::::::::::::::::::::::::::::::::::::::::::::::::::::::::::::::::::::::::::::::
%                                    STYLING                                     
%                        Alter Visual Appearance of Text                         
%................................................................................
\usepackage{microtype}		% Enhance micro-typography (e.g.: protrusion, expansion)
\usepackage{fancyhdr}		% Custom headers and footers
\pagestyle{plain}		% header: empty; footer: centered page number
\usepackage[dvipsnames]{xcolor}	% Access to 68 named colors
\usepackage{moresize}		% Additional font sizes
\usepackage{enumerate}		% Customize item style in 'enumerate'
\usepackage{multicol}		% Environment with multiple columns of text
\usepackage[normalem]{ulem}	% More underline styles (and aligns bottoms of multiple underlines)
				% 'normalem' prevents underlines from altering \emph
\usepackage{appendix}		% Provides an 'appendices' environment for easy formatting
\usepackage{romannum}		% Adds romam numbers


%::::::::::::::::::::::::::::::::::::::::::::::::::::::::::::::::::::::::::::::::
%                                     FLOATS                                     
%               Manage Floating Content (e.g.: figures and tables)               
%................................................................................
\usepackage{subcaption}	% Format multiple images(tables) in one figure(table) environment
			% NOTE: This package *automatically* loads the 'caption' package,
\captionsetup{
	labelfont = {small,bf},
	textfont  = {small,it}
}
%\subcaptionsetup{...}	% Uncomment to customize sub-captions separately

\usepackage{booktabs}	% Format tables following scientific convention
\usepackage{float}	% Provides [H] 'Here' float placement specifier (use with caution)


%:::::::::::::::::::::::::::::::::::::::::::::::::::::::::::::::::::::::::::::::::
%                                    GRAPHICS                                    
%                             Generate/Import Images                              
%.................................................................................
% NATIVE PROGRAMMATIC DIAGRAMS & PLOTS
\usepackage{tikz}		% Create vector graphics programmatically
\usetikzlibrary{arrows.meta}	% Additional arrow styles
\usepackage{pgfplots}		% Tool for creating function plots and charts
\pgfplotsset{compat=1.18}	% Enforce version 1.18 for consistency

% EXTERNAL IMAGE MANAGEMENT
\usepackage{graphicx}		% LaTeX standard for including external image files


% █▀▄▀█ ▄▀█ ▀█▀ █ █
% █ ▀ █ █▀█  █  █▀█
%:::::::::::::::::::::::::::::::::::::::::::::::::::::::::::::::::::::::::::::::::
%                                 MATH PACKAGES                                  
%.................................................................................
% ALGORITHMS
\usepackage[
	portuguese,
	ruled,
	vlined,
	linesnumbered
]{algorithm2e}	% Typeset formal algorithms
\SetKwInput{KwRequer}{Requer}
\SetKwInput{KwPrecondicao}{Pré-condição}
\SetKwInput{KwInicializa}{Inicializa}
\SetKwFor{Para}{para}{faça}{fim para}
\SetKwIF{Se}{SenaoSe}{Senao}{se}{então}{senão se}{senão}{fim se}
\DontPrintSemicolon

% NICHE MATH PACKAGES
\usepackage{xfrac}	% Slanted fractions (\sfrac)
\usepackage{cancel}	% Strike-through for symbols (\cancel and \cancelto)
\usepackage{bbm}	% Shorthand solution for blackboard boldfaced 1 (use \mathbbm{1})

%::::::::::::::::::::::::::::::::::::::::::::::::::::::::::::::::::::::::::::::::::
%                           CUSTOM THEOREM ENVIRONMENTS                            
%                                   (Color Coded)                                  
%..................................................................................
% 1. DEFINE COLOR CODE FOR EACH THEOREM-LIKE ENVIRONMENT (makes change of color choice effortless)
\colorlet{lmm}{DarkOrchid}	% For lemma env.
\colorlet{thm}{Tan}		% For theorem env.
\colorlet{prop}{Maroon}		% For proposition env.
\colorlet{crl}{RawSienna}	% For corollary env.
\colorlet{def}{OliveGreen}	% For definition env.
\colorlet{eg}{RoyalPurple}	% For example env.
\colorlet{ex}{RoyalPurple}	% For exercise env.
\colorlet{obs}{MidnightBlue}	% For remark env.
\colorlet{note}{PineGreen}	% For notation remark env.

%. . . . . . . . . . . . . . . . . . . . . . . . . . . . . . . . . . . . . . . . . 
% 2. DEFINE CUSTOM THEOREM STYLES
% 2.1. STYLE I - TAUTOLOGIES - FIRST NUMBER, THEN NAME
\newtheoremstyle{number-name}
{4mm}						% Space above env.
{4mm}						% Space below env.
{\itshape}					% Body font
{}						% Indent amount
{\bfseries}					% Head font
{:}						% Punctuation after head
{.3em}						% Space after head
{\thmnumber{#2}\,\thmname{#1}\thmnote{ (#3)}}	% Head spec

% 2.2. STYLE II - EXAMPLES & EXERCISES - FIRST NAME, THEN ITALIC NUMBER 
\newtheoremstyle{name-it_number}
{4mm}						% Space above env.
{4mm}						% Space below env.
{\itshape}					% Body font
{}						% Indent amount
{\bfseries}					% Head font
{}						% Punctuation after head
{.3em}						% Space after head
{\thmname{#1}\,{\itshape\thmnumber{#2}}}	% Head spec

% 2.3. STYLE III - REMARKS, SOLUTIONS & PROOFS - NO NUMBERING (should refer to a numbered env.)
\newtheoremstyle{no_number}
{4mm}						% Space above env.
{4mm}						% Space below env.
{}						% Body font
{}						% Indent amount
{\bfseries\itshape}				% Head font
{.}						% Punctuation after head
{.3em}						% Space after head
{}						% Head spec

% 2.4. STYLE IV - NOTATION REMARKS - NO NUMBERING (should refer to a numbered env.) AND ITALIC BODY 
\newtheoremstyle{no_number-it_body}
{4mm}						% Space above
{4mm}						% Space below
{\itshape}					% Body font
{}						% Indent amount
{\bfseries}					% Head font
{:}						% Punctuation after head
{.3em}						% Space after head
{}						% Head spec

%. . . . . . . . . . . . . . . . . . . . . . . . . . . . . . . . . . . . . . . . .
% 3. APPLY THEOREM STYLES AND DEFINE RESPECTIVE ENVIRONMENTS
% 3.1. STYLE I - TAUTOLOGIES - FIRST NUMBER, THEN NAME
\theoremstyle{number-name}
  \newtheorem{lemma}{\color{lmm}Lema}[section]
  \newtheorem{theorem}{\color{thm}Teorema}[section]
  \newtheorem{proposition}{\color{prop}Proposição}[section]
  \newtheorem{corollary}{\color{crl}Corolário}[section]
  \newtheorem{definition}{\color{def}Definição}[section]

% Manually force counters to share the 'lemma' counter (i.e.: shared counter for style I)
\makeatletter
\let\c@theorem\c@lemma
\let\c@proposition\c@lemma
\let\c@corollary\c@lemma
\let\c@definition\c@lemma
\makeatother

% 3.2. STYLE II - EXAMPLES & EXERCISES - FIRST NAME, THEN ITALIC NUMBER 
\theoremstyle{name-it_number}
  \newtheorem{exercise}{\textcolor{ex}{Exercício}}[section]
  \newtheorem{example}{\textcolor{eg}{Exemplo}}[section]

% Manually force 'example' to use the 'exercise' counter (i.e. shared counter for style II)
\makeatletter
\let\c@example\c@exercise
\makeatother

% 3.3. STYLE III - REMARKS, SOLUTIONS & PROOFS - NO NUMBERING
\theoremstyle{no_number}
  % Remark
  \newtheorem*{observation*}{\color{obs}Observação}
  % Solution
  \newtheorem*{solution}{\color{ex}Solução}
  % Proofs
  \newtheorem*{proofT}{\color{thm}Demonstração}
  \newtheorem*{proofL}{\color{lmm}Demonstração}
  \newtheorem*{proofC}{\color{crl}Demonstração}
  \newtheorem*{proofP}{\color{prop}Demonstração}
  \newtheorem*{proofE}{\color{eg}Demonstração}
  \newtheorem*{proofO}{\color{obs}Demonstração}

% 3.4. STYLE IV - NOTATION REMARKS - NO NUMBERING AND ITALIC BODY 
\theoremstyle{no_number-it_body}
  \newtheorem*{notation*}{\color{note}Notação}

%. . . . . . . . . . . . . . . . . . . . . . . . . . . . . . . . . . . . . . . . .
% EXTRA: CUSTOM QED SYMBOL (a small black square)
\renewcommand\qedsymbol{■}


%:::::::::::::::::::::::::::::::::::::::::::::::::::::::::::::::::::::::::::::::::
%                              CUSTOM MATH COMMANDS                               
%.................................................................................
\newcommand{\Bvec}[1]{\symbf{#1}}	% Bold symbol vector notation

% CONSTANTS
\newcommand{\e}{\mathscr{e}}		% Special notation for Euler's number
\newcommand{\im}{\mathscr{i}\,}		% Special notation for the imaginary unit

% RELATIONS
\newcommand{\longto}{\longrightarrow}	% Alias to note function domain/codomain (f: A --> B)

% Bellow, a colon-equals (and an analogous equals-colon) sign built 'from scratch', to be more compact and elegant
\newcommand*{\Defeq}{%
	\mathrel{%						% Wrapper so TeX treats/interprets the symbol as a relation
		\vcenter{%					% Vertically center the 'colon' part
			\baselineskip0.5ex%			% Manually set the vertical spacing between the dots
			\lineskiplimit0pt%			% Disable TeX's 'anti-collision' rule for full control
			\hbox{\scriptsize.}%			% The top dot
			\hbox{\scriptsize.}%			% The bottom dot
		}%
		{=}%						% The equals sign
	}%
}					% Compact colon-equals
\newcommand*{\Eqdef}{%
	\mathrel{%						% Wrapper so TeX treats/interprets the symbol as a relation
		{=}%						% The equals sign
		\vcenter{%					% Vertically center the 'colon' part
			\baselineskip0.5ex%			% Manually set the vertical spacing between the dots
			\lineskiplimit0pt%			% Disable TeX's 'anti-collision' rule for full control
			\hbox{\scriptsize.}%			% The top dot
			\hbox{\scriptsize.}%			% The bottom dot
		}%
	}%
}					% Compact equals-colon
\newcommand{\subnormal}{\triangleleft}	% Normal subgroup (left)
\newcommand{\supnormal}{\triangleright}	% Normal subgroup (right)

% UNARY OPERATIONS
\newcommand{\T}{{\mathsfup{T}}}				% Transpose symbol
\newcommand{\Tp}[1]{{#1}^{\mathsfup{T}}}	% Transpose operation

% LOGIC
\newcommand{\limplies}{\rightarrow}	% Conditional connector
\newcommand{\liff}{\leftrightarrow}	% Biconditional connector
\newcommand{\nand}{\uparrow}		% NAND connector
\newcommand{\nor}{\downarrow}		% NOR connector
\newcommand{\xor}{\veebar}		% XOR connector


% LOADED LAST TO AVOID DEPENDENCY ERRORS
%:::::::::::::::::::::::::::::::::::::::::::::::::::::::::::::::::::::::::::::::::
%                                  BIBLIOGRAPHY                                   
%                                Manage References                                
%.................................................................................
\usepackage[
	backend=biber,	% Use Biber backend (more powerful than BibTex)
	style=ieee,
	backref=true
]{biblatex}


%:::::::::::::::::::::::::::::::::::::::::::::::::::::::::::::::::::::::::::::::::
%                                  FINALIZATION                                   
%                 Packages That Patch Others. MUST BE LOADED LAST                 
%.................................................................................
\usepackage{hyperref}		% Add hyperlinks to references, ToC and URLs
\hypersetup{
	plainpages=false,	% Try to prevent duplicate page anchors
	colorlinks=true,	% true: colored links; false: boxed links
	linkcolor=NavyBlue,	% Internal links
	citecolor=NavyBlue,
	urlcolor=NavyBlue,
	breaklinks=true,
	pdfborder={0 0 0}
}

\usepackage[
	brazilian,
	nameinlink
]{cleveref}			% Objectively better than \autoref
% DEFINE 'cleveref' NAMES FOR CUSTOM THEOREM-LIKE ENVIRONMENTS
	% STYLE I ENVIRONMENTS
\crefname{lemma}{lema}{lemas}
\Crefname{lemma}{Lema}{Lemas}
\crefname{theorem}{teorema}{teoremas}
\Crefname{theorem}{Teorema}{Teoremas}
\crefname{proposition}{proposição}{proposições}
\Crefname{proposition}{Proposição}{Proposições}
\crefname{corollary}{corolário}{corolários}
\Crefname{corollary}{Corolário}{Corolários}
\crefname{definition}{definição}{definições}
\Crefname{definition}{Definição}{Definições}
	% STYLE II ENVIRONMENTS
\crefname{exercise}{exercício}{exercícios}
\Crefname{exercise}{Exercício}{Exercícios}
\crefname{example}{exemplo}{exemplos}
\Crefname{example}{Exemplo}{Exemplos}

\newcommand{\customref}[2]{%	% Simple alias for custom link text
	\hyperref[#2]{#1}%
}
